\begin{trapbox}
\textbf{\textcolor{CTrap}{误区一：把竖直差直接当作点到直线的距离}}

错误：写成 $PH=y_0-kx_0-b$。右侧可能为负，而且描述的是竖直方向，不一定垂直于斜线。
正确：先取 $R(x_0,kx_0+b)$，$PR=|y_0-kx_0-b|$；再由直角三角形换算
\[
 PH=\frac{|y_0-kx_0-b|}{\sqrt{1+k^2}}.
\]
理解页中的 $PR$ 与 $PH$ 已画在同一个三角形中：点不在直线上时，仅 $k=0$ 才有 $R=H$；点在直线上时则 $P=R=H$。

\medskip
\begin{minipage}{\linewidth}
\textbf{\textcolor{CTrap}{误区二：把对称公式的减号改成加号}}
\begin{center}
\begin{minipage}{0.98\linewidth}
\begin{tcolorbox}[colback=white,colframe=CStatement!25,boxrule=0.4pt,
  arc=1mm,left=3mm,right=3mm,top=3mm,bottom=3mm]
\centering
\textbf{把减号写成加号，会向远离对称轴的一侧走}\par\medskip
\begin{tikzpicture}[font=\small,>=Stealth]

\foreach \shift/\sgn in {0/1,6/-1} {
\begin{scope}[shift={(\shift,0)},x=1cm,y=.65cm]
\draw[CStatement,thick] (-1,0)--(1.5,0) node[right] {$y=0$};
\draw[CExample,thick,->] (0,\sgn)--(0,{-1*\sgn});
\draw[CTrap,dashed,thick,->] (0,\sgn)--(0,{3*\sgn});
\fill[CExample] (0,\sgn) circle (2pt) (0,{-1*\sgn}) circle (2pt);
\fill[CTrap] (0,{3*\sgn}) circle (2pt);
\node[left] at (0,\sgn) {$P(0,\sgn)$};
\node[left,text=CExample] at (0,{-1*\sgn}) {$P'(0,\pgfmathparse{-1*\sgn}\pgfmathprintnumber{\pgfmathresult})$};
\node[left,text=CTrap] at (0,{3*\sgn}) {$W(0,\pgfmathparse{3*\sgn}\pgfmathprintnumber{\pgfmathresult})$};
\end{scope}}

\end{tikzpicture}\par\medskip
同尺度取 $A=0,B=1,C=0$。实线是正确位移 $-2t(A,B)$，虚线是错误位移 $+2t(A,B)$。减号保证走向对称轴；实际向上还是向下由 $t$ 的正负决定。
\end{tcolorbox}
\end{minipage}
\end{center}

\end{minipage}
错误：将 $P'=P-2t(A,B)$ 写成 $P+2t(A,B)$，会朝错误方向走。
正确：保留减号，再根据 $t$ 的正负决定实际方向。它并不表示横纵坐标都一定减小。

\medskip
\textbf{\textcolor{CTrap}{误区三：把距离 $PH$、有向距离 $d$ 与比例 $t$ 混为一谈}}
\[
 PH=\frac{|Ax_0+By_0+C|}{\sqrt{A^2+B^2}},\qquad
 d=\frac{Ax_0+By_0+C}{\sqrt{A^2+B^2}},
\]
\[
 t=\frac{Ax_0+By_0+C}{A^2+B^2}
   =\frac{d}{\sqrt{A^2+B^2}}.
\]
错因：忘记 $(A,B)$ 未必长 $1$。正确：距离乘单位方向，或比例乘原方向向量，二者选一，不混用。

\medskip
\textbf{\textcolor{CTrap}{误区四：斜截式系数或符号读错}}

对 $y=kx+b$，一般式为 $kx-y+b=0$，故 $A=k,B=-1,C=b$。
若用 $\Delta=y_0-kx_0-b$、$s=\Delta/(1+k^2)$，则\textbf{$t=-s$}；
因此斜截式公式是 $P'=(x_0+2ks,y_0-2s)$，不能把 $s$ 不加区别地代入一般式中的 $t$。

\medskip
\textbf{\textcolor{CTrap}{误区五：中点在线上就认定对称}}

中点在直线上只保证“两侧代入值变号”，仍可能沿斜向连线。
正确：还须检查 $(x'-x_0,y'-y_0)$ 与 $(A,B)$ 平行。若点本就在直线上，则对称点就是它本身。
\end{trapbox}
